\clearpage
\section*{7. Kết luận và hướng phát triển}
\addcontentsline{toc}{section}{7. Kết luận và hướng phát triển}

\subsection*{7.1. Kết luận}

Bài thực hành đã xây dựng package ROS~2 độc lập
\texttt{tbe4\_face\_recognition}, kết hợp webcam, FER/TensorFlow Lite và
TurtleBot~4 qua giao diện topic chuẩn. Kiến trúc hai node giữ phần AI ở
virtualenv riêng và giữ vòng điều khiển 20~Hz không phụ thuộc thời gian
inference. Protocol JSON có session/sequence, filter kiểm tra đúng một khuôn
mặt và confidence lớn hơn 0,60, còn state machine one-shot ngăn hành động tự
lặp.

Năm profile chuyển động non-neutral dùng phản hồi odometry thay vì chỉ chạy
theo thời gian; neutral dừng robot. Watchdog observation/odometry, kiểm tra
sequence/session, pose jump, timeout, preview mặc định và zero Twist khi
cleanup tạo nhiều lớp dừng an toàn cho phạm vi bài thực hành.

Bằng chứng tự động cuối cho commit ứng dụng \texttt{c35bdd5} là 102 pytest
pass và 102 Colcon test với 0 lỗi, 0 thất bại, 0 bỏ qua. Webcam và FER/TFLite
đã chạy headless 12~giây trên frame thật mà không exception. Các kết quả này
chứng minh contract, thuật toán lọc, state machine, profile số và adapter hoạt
động theo test; chúng không thay thế nghiệm thu trực quan sáu biểu cảm và năm
quỹ đạo trên simulator.

\subsection*{7.2. Giới hạn}

\begin{itemize}[leftmargin=1cm,itemsep=0.2em]
  \item Phản ứng tức thời ở confidence lớn hơn 0,60 giảm độ trễ nhưng nhạy hơn
        với một lần phân loại sai so với thiết kế ổn định nhiều frame.
  \item FER2013 chỉ cung cấp nhóm nhãn biểu cảm tổng quát; độ chính xác phụ
        thuộc ánh sáng, góc mặt, che khuất và khác biệt giữa người dùng.
  \item Bài chỉ hỗ trợ một khuôn mặt. Nhiều mặt luôn dừng thay vì chọn một
        người điều khiển.
  \item Không có obstacle avoidance riêng và không chạy Nav2; mọi thử nghiệm
        motion cần vùng trống và giám sát trực tiếp.
  \item Chưa có đo lường thực nghiệm về accuracy, confusion matrix, latency
        dài hạn, sai số quãng đường/góc hoặc độ lặp lại trên robot thật.
  \item Cảnh báo phiên bản SciPy hệ thống/NumPy vẫn xuất hiện nhưng chưa làm
        hỏng đường inference đã thử.
\end{itemize}

\subsection*{7.3. Hướng phát triển}

Các bước tiếp theo phù hợp là thu một tập kiểm thử có sự đồng ý của người thực
hiện để đo confusion matrix; bổ sung hysteresis hoặc voting tùy chọn khi cần
giảm false trigger; đo latency camera--inference--Twist; đối chiếu odometry với
ground truth; và tích hợp safety layer/obstacle detection trước khi dùng trong
không gian có người. Với robot thật, cần thử giới hạn tốc độ thấp hơn, nút dừng
khẩn cấp và cơ chế arbitration cho \path{/cmd_vel}.

\subsection*{7.4. Liên kết}

\begin{itemize}[leftmargin=1cm]
  \item \textbf{Repository public:}
        \url{https://github.com/thanthe22019-oss/Tb4_face_recognition}
  \item \textbf{Bộ ảnh minh chứng:} \evidenceLink
  \item \textbf{Video demo:} \pending{bổ sung URL sau khi quay đủ sáu biểu cảm,
        năm profile và các tình huống dừng.}
\end{itemize}

\addcontentsline{toc}{section}{Tài liệu tham khảo}
\bibliographystyle{plain}
\bibliography{references}
